\chapter{Contest: Point-Level Instance Arbitration}
\label{ch:contest}

% PRESUPUESTO: ~8 pp. El capítulo más largo: es la contribución estrella.
% CONTRIBUCIÓN 4.
%
% ORDEN: este capítulo y el \ref{ch:geometric_consistency} son intercambiables.
% NO escribir referencias de orden ("como vimos en el capítulo anterior").
% NOTA: capítulo traducido al inglés (2026-07-29). Envuelto en otherlanguage{english}.

\begin{otherlanguage}{english}

\section{From instance to point}
\label{sec:contest_granularity}

As already seen, OVO separates two tasks: the \emph{geometric} reconstruction of the map, run by the
SLAM, and a \emph{semantic} layer of assignment and classification mounted on top of that
reconstruction. Each new keyframe contributes geometry, but the reconstruction does not accumulate
everything it sees without criterion: to keep the map coherent, before adding points it checks which
part of what it observed corresponds to already-reconstructed surface and discards it. Only regions
seen for the first time generate new points; what is already reconstructed is not duplicated, it
merely notes that it has been observed again. The result is a \emph{deduplicated} cloud: one point
per physical surface, not one per observation.

It is a sound design decision, for two reasons. It bounds the map's scale to the scene's area rather
than to the trajectory's length: a wall seen a hundred times is still a wall, not a hundred layers of
points. And it avoids stacking slightly misaligned reconstructions of the same surface, which would
blur it. That geometric permanence is well earned: a surface's position is a \emph{fact} that does
not change, and treating it as fixed is the right thing to do.

On top of that reconstruction the semantic layer is mounted. Each point carries an owner, the
instance it belongs to, and that owner is fixed by the first SAM mask that gathered enough votes to
claim the point, whether by creating a new instance or by adjudicating it to one that already existed
and under which it fell. The semantic layer thus inherits geometry's permanence: once assigned, the
point keeps its owner for the whole life of the map. But the two permanences are not worth the same.
A point's position is a measured fact; its owner is a hypothesis, that of a single mask at a single
instant, never revisited.

Hence the rigidity. However many keyframes later show that the correct owner is another, there is no
lever to change it. Dedup makes it worse: the surface already exists, so the mask that would later
claim it reprojects onto \emph{foreign} points and creates no points of its own to compete with. The
dispute never surfaces as geometry; it stays buried in the moment of the first assignment.

That is why the loop-closure fusion mechanism (\ref{sec:fm_scope_limit}) cannot reach these failures.
That \emph{merge} joins two instances when their clouds overlap, and dedup prevents that overlap from
existing: the misassigned points have no correct copy on top of them, they are the same points with
the wrong owner. With no two clouds to compare, the merge has nothing to trigger on. Three
situations, recurrent in practice, illustrate it (Figure~\ref{fig:contest_cases}).

\begin{figure}[htbp]
  \centering
  % PLACEHOLDER: sustituir cada \fbox por \includegraphics de un screenshot de Rerun (PNG),
  % cada uno un montaje Antes | Después sobre nube coloreada por instancia.
  % Escena + IDs de instancia por fijar (candidato para reusar como figura de fallo del Ch1).
  \newcommand{\caseph}[1]{\fbox{\parbox[c][2.2cm][c]{0.96\linewidth}{\centering\footnotesize\itshape #1}}}
  \begin{subfigure}{\linewidth}
    \centering\caseph{Before: contaminated chunk \quad$\vert$\quad After: chunk returned}
    \caption{Birth under a foreign mask.}
    \label{fig:contest_case_a}
  \end{subfigure}\\[4pt]
  \begin{subfigure}{\linewidth}
    \centering\caseph{Before: revisit split into two instances \quad$\vert$\quad After: reunited}
    \caption{Partial revisit.}
    \label{fig:contest_case_b}
  \end{subfigure}\\[4pt]
  \begin{subfigure}{\linewidth}
    \centering\caseph{Before: orphan fragment on the floor \quad$\vert$\quad After: reabsorbed}
    \caption{Inconsistent segmenter granularity.}
    \label{fig:contest_case_c}
  \end{subfigure}
  \caption{The three failure modes beyond the merge's reach (left), resolved by the contest (right).}
  \label{fig:contest_cases}
\end{figure}

\paragraph{Case A: birth under a foreign mask.}
An object enters the field of view and SAM segments it together with another that is not its own: a
single mask covers both. The points of the two objects are born under the same instance. Later, when
the object is seen unobstructed and SAM segments it well, the new instance is created correctly, but
the chunk trapped in the first one already has an owner and no one returns it. The contamination is
permanent (Figure~\ref{fig:contest_case_a}).

\paragraph{Case B: the partial revisit.}
An object is seen partially and leaves the scene before consolidating. On re-entering from another
side, the view is better and it is assigned a new instance that grows healthy. Even if in some later
keyframe SAM's masks segment well the chunk that had already been seen the first time, those points
still belong to the old instance. The two should be one, but they do not overlap (they are the same
surface seen at two moments, not two copies in the same place), so the merge does not consider them
candidates and never joins them (Figure~\ref{fig:contest_case_b}).

\paragraph{Case C: inconsistent segmenter granularity.}
SAM does not segment with a stable granularity: sometimes it splits hairs and breaks a large surface
(the floor is the typical case), other times it gives it whole. If the first sighting of an area
over-splits it and the rest of the trajectory sees it complete, the initial fragment remains a tiny
instance \emph{inside} the large one. It falls within it in every subsequent frame, but having an
owner already it is not reabsorbed: the merge, again, sees no duplicates to join
(Figure~\ref{fig:contest_case_c}).

The three share a root, and it is not a poorly calibrated heuristic: they are failures \emph{outside
the domain} of the existing mechanisms. Each one fulfills its responsibility, dedup the geometric
coherence and loop-closure fusion the duplications from drift or revisit, but none \emph{refines} the
assignments in light of what is observed afterward. That is the gap the contest fills, and it does
not go through undoing dedup but through making a point's owner revisable by evidence instead of fixed
by the first mask: the geometry stays frozen, the semantics stop being so. The refinement rests on
two decisions.

\begin{enumerate}
  \item \textbf{Decision granularity.}
  A partial failure is only fixed by a partial decision. If the unit of decision is the whole
  instance, \enquote{this chunk belongs to B} cannot even be posed: either the whole instance moves
  or none of it. By lowering the decision to the point, each point can change owner on its own, and
  the classic \emph{merge} becomes the extreme case in which it turns out that all the points of an
  instance migrate to another.

  \item \textbf{Criterion to revoke.}
  Precisely because the first decision may be wrong, revoking it lightly, on the evidence of a single
  frame, reintroduces the same problem in reverse. The only reliable basis for moving a point's owner
  is \emph{accumulated evidence}: observing across many keyframes which instance consistently claims
  that point, and only then deciding. From blindly trusting each mask we move to trusting the history
  of masks.
\end{enumerate}

\medskip
\noindent\emph{Precondition.} All of this presupposes stable geometry and reliable tracking. Evidence
accumulates by reprojecting masks onto the points, so if the map drifts or the tracking fails, that
reprojection lands where it should not and generates disputes that answer to no real conflict: a
point's history stops being comparable with itself and the reasoning breaks. But this is, again,
someone else's responsibility. Keeping the geometry healthy beneath the points is the task of the
geometric refresh and submaps (\ref{ch:geometric_consistency}); the contest takes that condition as
solved and deals only with its own, arbitrating the dispute.

\section{The contest mechanism}
\label{sec:contest_mechanism}

The contest is the mechanism that assumes the missing responsibility: it registers the conflicts
between instances, accumulates the evidence of each, and, from that evidence, decides what to do. It
is all organized around a pair of roles. When two instances dispute a point, the one that currently
owns it is the \emph{defender} $D$ and the one claiming it is the \emph{challenger} $C$. The unit
over which the mechanism reasons and decides is not a lone instance but the ordered pair $(D, C)$:
all the machinery (the evidence it stores, the signals it computes, the decision it emits) is
referred to that pair. Its output is a per-pair \emph{verdict}, $v(D, C) \in \{\textsc{merge},
\textsc{split}, \textsc{no-action}\}$, that translates the accumulated evidence into a concrete
action on the map or into its absence.

The mechanism breaks down into three stages, each responsible for one thing and no more. The
\emph{registry} (\ref{sec:contest_registry}) records the dispute point by point and aggregates it per
pair, without deciding. The \emph{discriminator} (\ref{sec:contest_discriminator}) reads that
evidence and emits the verdict, without touching the map. The \emph{actuator}
(\ref{sec:contest_actuator}) executes the verdict on the map. Registry and discriminator are pure,
and that purity is deliberate: the dispute is counted and judged separately from where it is applied,
and only the last stage modifies anything (Figure~\ref{fig:contest_pipeline}).

\begin{figure}[ht]
  \centering
  % PLACEHOLDER: diagrama del pipeline del mecanismo, a dibujar (TikZ o drawio -> PDF).
  \fbox{\parbox[c][3.2cm][c]{0.9\linewidth}{\centering\footnotesize\itshape
  Contest pipeline: tracking $\to$ registry (per-pair signals) $\to$ discriminator (verdict) $\to$
  actuator (mutated map). Registry and discriminator pure; only the actuator touches the map.}}
  \caption{The three stages of the contest and the data flow between them.}
  \label{fig:contest_pipeline}
\end{figure}

\subsection{The registry}
\label{sec:contest_registry}

The registry is \emph{online}: it is built at the same time as the map, with no later pass. At each
tracking step, each SAM mask is reprojected onto the visible points and resolves to an instance, the
majority one among those already owned under the mask. Let $p$ be a point owned by $D = o(p)$; if it
falls under a mask that resolves to $C \neq D$, there is a \emph{grab} of $p$ by $C$ against $D$
(Figure~\ref{fig:contest_grab}). It is not applied: it is recorded, adding a grab to $C$'s count over
that point.

\begin{figure}[ht]
  \centering
  % PLACEHOLDER: diagrama esquemático del robo, a dibujar (TikZ o drawio -> PDF).
  \fbox{\parbox[c][2cm][c]{0.85\linewidth}{\centering\footnotesize\itshape
  Grab diagram: SAM mask reprojected onto the cloud; a point whose defender differs from the mask's
  challenger; the grab is recorded.}}
  \caption{The grab: the contest's atomic observation.}
  \label{fig:contest_grab}
\end{figure}

The registry has two levels. The first is \emph{per point}: three counters indexed by each point's
permanent id (Table~\ref{tab:contest_primary}), the base truth that drips step by step. Each point
carries its count of grabs (per challenger), of claims, and of sightings, and nothing else is stored
at this level; everything else is derived from here.

\begin{table}[ht]
  \centering
  \small
  \begin{threeparttable}
  \renewcommand{\arraystretch}{1.5}
  \begin{tabular}{@{}l l p{7.4cm}@{}}
    \toprule
    \textbf{Signal} & \textbf{Store} & \textbf{Steps counted} \\
    \midrule
    \textbf{\texttt{grabs}}     & $\texttt{point\_id}\mapsto\{C\!:\#\mathrm{KF}\tnote{*}\}$ &
      The point, owned by $D$, fell under $C$'s mask (a grab) \\
    \textbf{\texttt{claims}}    & $\texttt{point\_id}\mapsto\#\mathrm{KF}$ &
      The point fell under any used mask ($C$ or $D$) \\
    \textbf{\texttt{sightings}} & $\texttt{point\_id}\mapsto\#\mathrm{KF}$ &
      The point reprojected into view, with or without a mask \\
    \bottomrule
  \end{tabular}
  \begin{tablenotes}
    \footnotesize\item[*]\itshape $\#\mathrm{KF}$: number of
      keyframe steps.
  \end{tablenotes}
  \end{threeparttable}
  \caption{Primary signals: per-point counters accumulated in the store.}
  \label{tab:contest_primary}
\end{table}

The second level aggregates those counters \emph{per pair} $(D, C)$, the unit of the decision. But
before aggregating, signal must be separated from noise, and that is the job of \textbf{firmness}: a
point's grab counts as firm only if the challenger took it over enough steps (\emph{evidence}, a
minimum number of grabs) and those grabs are a real fraction of the point's claims
(\emph{commitment}, the grabs/claims ratio above a threshold). A glancing brush, or a flicker of a
few steps, is not firm. A pair's number of firm points is \texttt{firm\_points}, and on that subset
all the secondary signals are built (Table~\ref{tab:contest_secondary}). These live cached per
defender: a grab only dirties its defender, whose signals are recomputed at the next refresh, without
re-deriving the rest.

\begin{table}[ht]
  \centering
  \small
  \begin{threeparttable}
  \renewcommand{\arraystretch}{1.5}
  \begin{tabular}{@{}l p{11.6cm}@{}}
    \toprule
    \textbf{Signal} & \textbf{How it is computed} \\
    \midrule
    \textbf{\texttt{firm\_points}} & points of $D$ that $C$ steals firmly \\
    \textbf{\texttt{total\_grabs}} & sum of all $C$'s grabs over $D$'s points (firm or not) \\
    \textbf{\texttt{containment}}  & $\mathrm{firm\_points}\,/\,|D|\tnote{*}$ \\
    \textbf{\texttt{confidence}}   & per-point $\mathrm{grabs}_C/\mathrm{claims}$, averaged over the disputed points \\
    \textbf{\texttt{exclusivity}}  & firm points disputed only by $C$ $\,/\,\mathrm{firm\_points}$ \\
    \bottomrule
  \end{tabular}
  \begin{tablenotes}
    \footnotesize\item[*]\itshape $|D|$: number of points of instance $D$.
  \end{tablenotes}
  \end{threeparttable}
  \caption{Secondary signals: per-pair $(D,C)$ aggregates over firm points.}
  \label{tab:contest_secondary}
\end{table}

This is the whole contest registry: a per-point stream of grabs that drips at each step, three
primary per-point counters, and a per-pair aggregate cached per defender. At no point has the map
been touched or an assignment revoked; the registry is a faithful accounting of the dispute,
deliberately separated from the decision. Turning that evidence into a verdict is the discriminator's
job.

\subsection{The discriminator}
\label{sec:contest_discriminator}

The discriminator takes a defender's pairs $(D, C)$, with their signals already computed, and emits
the verdict. It is a pure function: it looks neither at the map nor at the descriptor, only at the
evidence the registry accumulated. All of it is organized around a single quantity, the containment
$\mathrm{containment}(D, C)$: what fraction of the defender the challenger firmly steals. Containment
decides the \emph{shape} of the action, and confidence and exclusivity decide whether there is
\emph{enough evidence} to execute it.

Before reasoning anything, a mass filter discards noise: a pair with fewer than $\tau_{\text{mass}}$
raw grabs ($\mathrm{total\_grabs} < 50$) is not a dispute, it is a brush, and is ignored. Over the
pairs that survive, containment splits the decision into three bands.

\begin{enumerate}
  \item \textbf{High containment} ($\mathrm{containment} \geq 0.6$): the challenger takes almost the
  whole defender, so the defender is a \emph{fragment} of it (cases A and C). The natural action is
  \textsc{merge}, but only if there is a reliable owner: a confidence $\mathrm{confidence} \geq 0.5$
  is required (on average, each point is stolen in the majority of its claims) and that all firm
  challengers are, via fusion, a single instance. If several distinct objects contain it there is no
  clear owner and \textsc{no-action} is emitted.

  \item \textbf{Partial containment} ($0.4 \leq \mathrm{containment} < 0.6$): the challenger takes a
  large part, but not so much that magnitude alone suffices. Here the threshold rises: \textsc{merge}
  only if the chunk is \emph{clean}, that is with high confidence ($\mathrm{confidence} \geq 0.6$)
  and exclusive ($\mathrm{exclusivity} \geq 0.8$, almost all its points are disputed by this
  challenger and no one else), and again with a single root. The magnitude evidence is weaker, so it
  is compensated by demanding more cleanliness in the other two signals.

  \item \textbf{Low containment} ($\mathrm{containment} < 0.4$): the defender \emph{retains} most of
  its points but cedes exclusive chunks to one or several neighbors. It is no one's fragment, it is
  an object dragging pieces of others (case B). The action is \textsc{split}, and not one but as many
  as needed: for each challenger with a chunk of high confidence ($\mathrm{confidence} \geq 0.6$) and
  very exclusive ($\mathrm{exclusivity} \geq 0.9$) that chunk is detached toward it. If no challenger
  qualifies, the pair is on the edge and \textsc{no-action} is emitted.
\end{enumerate}

The exclusivity threshold rises from $0.8$ in the partial \textsc{merge} to $0.9$ in the
\textsc{split}, and not on a whim: a split cuts geometry in two and is more expensive to undo if
unwarranted, so the chunk's boundary is required to be sharper before cutting. The floor
$\mathrm{containment} = 0.4$ separating split from merge is empirical, fixed on the tuning set
(\ref{sec:contest_results}).

Seen as a whole, the classic \textsc{merge} reappears as the extreme of this scale: the case in which
containment is total and the defender migrates whole. The discriminator merely extends that same idea
to intermediate and low containments, where the decision stops being to move the whole instance and
becomes to move exactly the points the evidence supports.

\subsection{The actuator}
\label{sec:contest_actuator}
% Ejecuta el veredicto sobre el mapa: MERGE (union-find de instancias) y SPLIT (mueve los
% split_points al challenger). Última etapa, la única que muta. Aquí se cuela el timing/coste
% real (Real-time ya no es sección propia): coste por refresco, cuello de botella y rediseño
% online. Doc: docs/contest_timing_analysis.html

\section{Results}
\label{sec:contest_results}

The discriminator's thresholds are empirical, so all the numbers are reported on a set of scenes the
tuning never saw. Replica's eight scenes are split into tuning (office0, office3, office4, room2),
the only ones on which the thresholds are fixed, and held-out (office1, office2, room0, room1), which
stay closed throughout the tuning and are opened only once, at the end. No threshold is touched by
looking at the held-out; the $\mathrm{AP}$ that is defended is always the held-out's, so it measures
generalization and not overfitting.

% PRELIMINAR: números a recalcular antes de la entrega (placeholder para ver la maqueta).
On Replica's clean trajectory, with no injected jumps, the contest raises the global
$\mathrm{AP}_{\text{agnostic}}$ from $0.234$ to $0.249$ ($+0.015$, $+6.4\%$) over the raw map
(\emph{baseline} with no fusion on simulated GT). The improvement appears in all eight scenes, and
holds on the held-out: the mean gain is $+9.7\%$ on the tuning scenes and $+4.55\%$ on the held-out
ones, positive in both cases (Table~\ref{tab:contest_ap_scenes}).

\begin{table}[ht]
  \centering
  \small
  \renewcommand{\arraystretch}{1.2}
  \begin{tabular}{@{}l r r r r@{}}
    \toprule
    \textbf{Scene} & \textbf{Baseline} & \textbf{Contest} & \textbf{$\Delta$ abs} &
      \textbf{$\Delta$ \%} \\
    \midrule
    \multicolumn{5}{@{}l}{\textbf{Tuning}} \\
    office0 & 0.199 & 0.220 & $+0.021$ & $+10.6\%$ \\
    office3 & 0.204 & 0.210 & $+0.006$ & $+2.9\%$  \\
    office4 & 0.236 & 0.280 & $+0.044$ & $+18.6\%$ \\
    room2   & 0.232 & 0.247 & $+0.015$ & $+6.5\%$  \\
    \addlinespace[1pt]
    \textbf{mean} & \textbf{0.218} & \textbf{0.239} & $\mathbf{+0.021}$ & $\mathbf{+9.7\%}$ \\
    \midrule
    \multicolumn{5}{@{}l}{\textbf{Held-out}} \\
    office1 & 0.170 & 0.175 & $+0.005$ & $+2.9\%$  \\
    office2 & 0.210 & 0.225 & $+0.015$ & $+7.1\%$  \\
    room0   & 0.269 & 0.287 & $+0.018$ & $+6.7\%$  \\
    room1   & 0.344 & 0.349 & $+0.005$ & $+1.5\%$  \\
    \addlinespace[1pt]
    \textbf{mean} & \textbf{0.248} & \textbf{0.259} & $\mathbf{+0.011}$ & $\mathbf{+4.55\%}$ \\
    \midrule
    \textbf{All} & \textbf{0.234} & \textbf{0.249} & $\mathbf{+0.015}$ & $\mathbf{+6.4\%}$ \\
    \bottomrule
  \end{tabular}
  \caption{$\mathrm{AP}_{\text{agnostic}}$ per Replica scene, baseline vs.\ contest.}
  \label{tab:contest_ap_scenes}
\end{table}

The \emph{class-aware} metrics (mIoU / mAcc) do not worsen: they move slightly upward (mIoU $+0.2$,
mAcc $+0.3$ globally, Table~\ref{tab:contest_semantic}). This is consistent with the nature of the
mechanism: the contest reorganizes point ownership among instances, it does not change the semantic
label each one receives, so its effect lives in the granularity of the instances (what the agnostic
$\mathrm{AP}$ measures) and leaves the per-class mIoU practically intact.

\begin{table}[ht]
  \centering
  \small
  \renewcommand{\arraystretch}{1.2}
  \begin{tabular}{@{}l r r r r r r@{}}
    \toprule
    & \multicolumn{3}{c}{\textbf{mIoU}} & \multicolumn{3}{c}{\textbf{mAcc}} \\
    \cmidrule(lr){2-4} \cmidrule(lr){5-7}
    \textbf{Scene} & Base & Contest & $\Delta$ & Base & Contest & $\Delta$ \\
    \midrule
    office0 & 25.4 & 25.9 & $+0.5$ & 33.9 & 34.3 & $+0.4$ \\
    office1 & 16.2 & 16.1 & $-0.1$ & 31.7 & 31.7 & $+0.0$ \\
    office2 & 24.6 & 25.9 & $+1.3$ & 32.4 & 34.0 & $+1.5$ \\
    office3 & 25.8 & 25.9 & $+0.0$ & 32.3 & 32.3 & $+0.0$ \\
    office4 & 43.1 & 42.7 & $-0.3$ & 55.4 & 55.8 & $+0.4$ \\
    room0   & 33.2 & 33.6 & $+0.4$ & 39.9 & 40.3 & $+0.4$ \\
    room1   & 38.6 & 40.2 & $+1.6$ & 53.8 & 56.1 & $+2.3$ \\
    room2   & 26.6 & 26.6 & $+0.0$ & 35.1 & 35.1 & $+0.0$ \\
    \midrule
    \textbf{All} & \textbf{27.1} & \textbf{27.3} & $\mathbf{+0.2}$ & \textbf{40.2} & \textbf{40.5} & $\mathbf{+0.3}$ \\
    \bottomrule
  \end{tabular}
  \caption{Class-aware semantic metrics (\%) per Replica scene, baseline vs.\ contest.}
  \label{tab:contest_semantic}
\end{table}

\end{otherlanguage}
% Validación en AISLAMIENTO, sobre trayectoria sana y sin saltos: se atribuye la mejora
% al mecanismo, no a su interacción con el ruido. COROLARIO, que hay que escribir: el
% contest PRESUPONE geometría sana y tracking fiable -> el sustrato no es un preámbulo,
% es una PRECONDICIÓN.
%
% ABLACIÓN 2x2 (el experimento estrella de la memoria), sobre saltos inyectados:
%
%                    | sin submapas            | con submapas
%   .................+.........................+...........................
%   sin contest      | baseline degradado      | ¿el submapa solo, ayuda?
%   con contest      | contest sobre sustrato  | la propuesta completa
%                    | roto -> debería ir mal  |
%
% La celda de abajo-izquierda es la que DEMUESTRA que el orden causal es real y no
% retórica. Asegurar este experimento sí o sí.
